% ============================================================
%  Section 7: Evaluation Metrics
%  H&E to Sirius Red Virtual Staining Paper
% ============================================================

\subsection{Evaluation Metrics}
\label{sec:metrics}

Since the ultimate objective of virtual staining is to support downstream biological
analysis, task-specific evaluation provides the most direct measure of model performance.
A translation that faithfully preserves the relevant tissue signal will produce
biologically consistent quantitative readouts when the same analysis pipeline is applied
to both real and virtual images.
We supplement this with structural, perceptual, and distributional measures that
characterise the visual quality of the generated stains.

Because H\&E and SR sections are non-adjacent cuts (Sec.~\ref{sec:dataset}), direct
pixel-level comparison cannot serve as ground truth.
We therefore co-register the test-set H\&E and SR WSIs with the VALIS feature-based
framework~\cite{gatenbee2023virtual}, which reaches tile-scale but not pixel-exact
correspondence, so residual misalignment inflates apparent error for any pixel-level
comparison metric.

\noindent\textbf{Perceptual and Distributional Metrics:} Given these constraints, we report patch-based SSIM~\cite{wang2004ssim},
LPIPS~\cite{zhang2018lpips}, and FID~\cite{heusel2017fid}; each captures a different
aspect of perceptual or distributional quality and each has known limitations for
virtual staining.
Patch-based SSIM computes structural similarity over local patches to reduce
sensitivity to residual misregistration, and LPIPS measures VGG-based deep-feature
distance over the co-registered tile pairs without strict pixel pairing.
FID compares real and generated SR distributions using InceptionV3 features, capturing
population-level realism but unable to detect localised semantic failure on individual
tiles.
Because these InceptionV3 features are trained on natural images rather than histology,
FID can reward distributional realism while missing collagen-specific structure; we
therefore treat CPA~MAE, not FID, as the primary metric.

\noindent\textbf{Task-Specific Evaluation: Collagen Proportionate Area (CPA) Agreement.}
Although the two sections cannot be compared pixel-by-pixel, they share the same
underlying collagen content and therefore allow a paired comparison at the level of
biological quantity. To realise this, a frozen %nnU-Net~v2~\cite{isensee2021nnunet}
collagen segmentation model~\cite{deng2026} trained on real SR WSIs is applied
identically to (i)~the real SR test WSIs and (ii)~the virtual SR WSIs produced by each
model configuration.
Deep learning segmentation models of this kind reach Dice scores of $0.94$
on Sirius Red~\cite{deng2026}, demonstrating that automated CPA measurement
is clinically viable.
For each WSI, CPA is computed as the fraction of tissue pixels labelled as SR-positive
collagen (background excluded).
CPA is the standard quantitative readout for hepatic fibrosis assessment and is
directly comparable across slides~\cite{huang2013sirius}.

The primary task-specific metric is the \textit{mean absolute error}~(MAE) between
paired real-SR and virtual-SR CPA values:
\begin{equation}
  \mathrm{MAE}_{\mathrm{CPA}} = \frac{1}{N}\sum_{i=1}^{N}
  \bigl|\mathrm{CPA}^{\mathrm{virtual}}_{i} - \mathrm{CPA}^{\mathrm{real}}_{i}\bigr|,
  \label{eq:cpa_mae}
\end{equation}
where $N$ is the number of held-out paired specimens.
A lower CPA~MAE indicates that the virtual stain supports collagen quantification
consistent with the real stain, serving as the primary biologically grounded benchmark
for model evaluation in this study.

\noindent\textbf{Expert Reader Study.}
Automated metrics are complemented by a blinded reader study with three expert readers,
two clinical liver experts and one reader experienced with histological images, conducted on the
held-out set for the study-best and study-worst configurations.
Readers first decided whether each SR tile was real or generated, then rated whether the
collagen in a generated tile is biologically plausible given its matched H\&E input on a
forced-choice four-point scale.
The full protocol is given in Supplementary Sec.~6.


% ============================================================
%  New bibliography entries for this section
%  (merge with references.bib)
% ============================================================
%
% @inproceedings{heusel2017fid,
%   author    = {Heusel, Martin and Ramsauer, Hubert and Unterthiner, Thomas and Nessler, Bernhard and Hochreiter, Sepp},
%   title     = {{GANs} Trained by a Two Time-Scale Update Rule Converge to a Local {Nash} Equilibrium},
%   booktitle = {NeurIPS},
%   year      = {2017},
%   eprint    = {1706.08500}
% }
%
% @inproceedings{szegedy2016inception,
%   author    = {Szegedy, Christian and Vanhoucke, Vincent and Ioffe, Sergey and Shlens, Jon and Wojna, Zbigniew},
%   title     = {Rethinking the Inception Architecture for Computer Vision},
%   booktitle = {CVPR},
%   year      = {2016}
% }
%
% @inproceedings{zhang2018lpips,
%   author    = {Zhang, Richard and Isola, Phillip and Efros, Alexei A. and Shechtman, Eli and Wang, Oliver},
%   title     = {The Unreasonable Effectiveness of Deep Features as a Perceptual Metric},
%   booktitle = {CVPR},
%   year      = {2018},
%   eprint    = {1801.03924}
% }
%
% @article{simonyan2015vgg,
%   author    = {Simonyan, Karen and Zisserman, Andrew},
%   title     = {Very Deep Convolutional Networks for Large-Scale Image Recognition},
%   booktitle = {ICLR},
%   year      = {2015},
%   eprint    = {1409.1556}
% }
%
% @article{wang2004ssim,
%   author    = {Wang, Zhou and Bovik, Alan C. and Sheikh, Hamid R. and Simoncelli, Eero P.},
%   title     = {Image Quality Assessment: From Error Visibility to Structural Similarity},
%   journal   = {IEEE Transactions on Image Processing},
%   volume    = {13},
%   number    = {4},
%   pages     = {600--612},
%   year      = {2004},
%   doi       = {10.1109/TIP.2003.819861}
% }
%
% @article{isensee2021nnunet,
%   author    = {Isensee, Fabian and Jaeger, Paul F. and Kohl, Simon A. A. and
%                Petersen, Jens and Maier-Hein, Klaus H.},
%   title     = {nnU-Net: a self-configuring method for deep learning-based
%                biomedical image segmentation},
%   journal   = {Nature Methods},
%   volume    = {18},
%   pages     = {203--211},
%   year      = {2021},
%   doi       = {10.1038/s41592-020-01008-z}
% }
%
% (Shared entries already listed:
%  huang2013sirius)
